% The long form of §8, as it stood on 2026-11-10 before the section was reorganized by direction
% (receive, send, forwarding; full speed, against message size, on one core and on several: the
% owner, 2026-11-10; outline.md "§8 Evaluation, by direction"). NOT \input by main.tex.
%
% It is kept because it holds, with their CLAIMS rows, the paragraphs the short form leaves out:
% the eight methodology traps, the paced cells, the full transmit table (tab:tx), the send path's
% mechanism rows (S7-S9, S13-S25), the splice series (N1-N3) and compaction's history (F3). To
% bring a paragraph back, move it into 08-evaluation.tex; its figure (figures/tx.tex with the
% paced panel) is at git revision 1f81527.
%
% ---------------------------------------------------------------------------------------------
% §8, 3.25 pages, in the order of the title (outline.md "§8 Evaluation, in the order of the
% title"). Every number is a CLAIMS.md row, quoted as written, with its status: pinned and
% measured rows are results; in-placement rows say so; argued rows say "by design"; pending and
% unmeasured rows are \todom. Row ids are in comments.
\section{Evaluation}
\label{sec:eval}

The evaluation follows the title. \S\ref{sec:eval:method} says how we measure, because measuring
zero copy honestly turned out to need care. \S\ref{sec:eval:ubiquitous} asks where each
mechanism gives zero copy at all. \S\ref{sec:eval:tx} measures the absent overhead on the send
path, where every mechanism can be compared on one machine. \S\ref{sec:eval:fwd},
\S\ref{sec:eval:retention} and \S\ref{sec:eval:store} measure what lent memory can do:
forwarding, holding and fan-out, and the store.

\subsection{Measuring zero copy honestly}
\label{sec:eval:method}

\paragraph{Testbed.} Three virtual machines in a public cloud, each with two virtio-net
interfaces and a 25\,Gbps allocation per interface, form a triangle: a source, the machine
under test, and a sink, with the two data legs on separate interfaces. The machine under test
is a two-socket guest; its data-path threads are pinned, and the NIC's interrupts are pinned
to the same socket. The kernels are the evaluation kernel of \S\ref{sec:impl}, on Linux
7.3-rc6; the source runs a stock kernel. \todov{machine shapes, vCPU counts and link rates at submission}

\paragraph{Metric.} The cost we report is the whole host's CPU time per gigabyte delivered,
measured as total minus idle over all CPUs of the machine under test, with the data-path
process's own share alongside. Process accounting alone makes kernel-resident mechanisms look
free, because softirq work lands on whichever CPU takes the interrupt and sockmap's data path
runs entirely in softirq and a kernel worker. Throughput is bytes delivered at the sink; paced
runs hold a fixed rate so that mechanisms are compared at equal offered load, since unpaced
runs flatter mechanisms with large per-operation costs. Every cell is three runs.

\paragraph{Eight traps.} Each would have changed a conclusion, and each is a methodology point
for anyone measuring zero copy on Linux. (1)~\emph{Interrupt time booked as idle.} On the 7.3
kernels, interrupt and softirq work on a tickless idle CPU counts as idle unless interrupt-time
accounting is enabled, and the idle loop over-counts idle by part of a tick per sleep, so
``total minus idle'' under-reports every receive-side cost. We enable interrupt-time
accounting, measure idle from idle entry to idle exit, and validate the accounting on each
kernel with a pinned thread of known runtime and a NIC interrupt CPU of known load before any
campaign. (2)~\emph{The packet tap.} A DHCP client on every stock cloud image holds a packet
socket on each interface, and a zero-copy packet must be copied before it is cloned for a tap:
every \sendzc on these machines was silently a copy until the tap was taken out of the transmit
path. With the interrupts pinned, the tap costs \sendzc 15\% at 1\KiB and 32\% at 4\KiB at
0.5\GBs, 41\% at 1\GBs with 64\KiB, and 60--149\% at full speed; it costs the named mode
1--5\%. % CLAIMS F7, pinned
(3)~\emph{The usage report.} \iouring's report of which zero-copy sends were copied over-counts
copies when several sends share a packet, so we count real copies with a probe on the kernel's
one copy routine instead. (4)~\emph{Interleaving.} Several zero-copy sends in flight on one
socket interleave their bytes under backpressure, because the kernel retries a partial send by
itself; a stream verified byte by byte exposed it. Every sender here keeps order with linked
requests or one request in flight, and every run verifies the stream at the sink.
(5)~\emph{Interrupt placement.} Where a host's NIC interrupts land, chosen by the interrupt
balancer and RSS at random per run, moves CPU per gigabyte by 10--25\% for every sender, which
is larger than several effects reported in the literature; the campaigns below pin the
interrupts, and earlier campaigns are read within one placement. Pinning them on the wrong node
is worse than leaving them to chance: with every interrupt on the remote NUMA node, every
zero-copy sender costs 20--52\% more per gigabyte and the copying sender 33--97\% more, while
the differences between senders stay in the same range. % CLAIMS M1, M2, M3
(6)~\emph{Withheld notifications.} At 1\KiB paced, \iouring's notification stacking withholds
\sendzc's buffer-reuse notifications for most of a run: each send joins the notification chain
of the unsent tail packet, and the chain completes only when the write queue drains, so most
runs deliver no notification in the measurement window and the rest deliver over a hundred
thousand at once and overflow the completion queue. A paced 1\KiB \sendzc baseline is therefore
not steady unless the chains are bounded, and every \sendzc sender below bounds them with
\code{MSG\_EOR} every megabyte. Bounding also cut the non-vectored \sendzc senders' cost at
0.5\GBs, by 17\% and 18\% for the unregistered sender at 1 and 4\KiB and by 6\% and 17\% for
the registered one, all of it in the sending process, although at 4\KiB nothing was withheld:
long chains cost by themselves. The named mode has no notification to bound.
% CLAIMS S16, measured
(7)~\emph{The sender's slow state.} The sender's vCPU enters, at random, a state in which
everything on it costs 15--30\% more, for every kernel and sender: memory latency is 12--27\%
higher and VM exits 19--35\% more frequent, while register loops are unchanged. It held 6 of
36, 11 of 35 and 6 of 36 runs on three kernels at 1\KiB paced. Paced cells are classified by
the sender CPU's level and read in the fast state with the slow runs counted; the paced numbers
below are means over all runs, and where leaving the slow runs out moves one by more than a
point we say so. % CLAIMS M5, pinned
(8)~\emph{The copy's two costs.} On the receiving machine, \code{recv()}'s copy has two costs,
drawn per run and set by nothing the guest controls: the same kernel, boot and receive queue
give both. The dear copy costs 0.103--0.104 CPU-seconds per gigabyte more in the receiving
process and 0.118--0.141 on the host; it held 11 of 18 paced runs on one kernel, with none
between the two, and 14 of 18 on another, with one between. Only the receiver that reads the
payload splits: next to a cheap and a dear \code{recv()} run alike, the other receivers cost
0.96--1.05 of their own medians. We infer, without having measured it, that the host's
placement of the data sets the cost. Every comparison against \code{recv()} below says which
copy it is against. % CLAIMS M6, measured

\subsection{Ubiquitous}
\label{sec:eval:ubiquitous}

\paragraph{Where zero copy registers.} \todom{the availability table across physical and
virtual drivers, VMs of three clouds, loopback and veth: whether zero-copy receive registers,
whether the mapped mode registers, and the share of bytes the named mode captures by reference}
On the virtio-net interfaces of all three machines, with the zero-copy receive code built in,
registering a receive queue is refused with \code{EOPNOTSUPP}: the driver has no header/data
split and no flow steering, and the refusal is the device's, since a nonexistent interface
returns \code{ENODEV}. The named mode runs on the same interfaces unchanged, and every result in
this paper was measured there. % CLAIMS F5, measured

\paragraph{Receive on virtio-net.} Four receivers take the same streams on the receiving
machine: \code{recv()}; a named sink that discards what arrives, so that bytes are dropped in
the kernel and never reach the process; named capture, which takes each chunk as an object and
frees it; and the mapped mode, one multishot zero-copy receive with slots returned through the
refill ring. With a copying sender on the other machine, the receivers interleaved in time
and six runs per cell over three boots, the discarding sink costs the receiving host
0.285--0.302 CPU-seconds per gigabyte at 0.5 and 1\GBs. \code{recv()} costs 0.471--0.492 in
the eleven runs whose copy was dear (trap 8), 1.61--1.65$\times$ the sink's cost and
2.36--2.49$\times$ its cost in the receiving process, and 0.330--0.374 in the seven whose copy
was cheap, 1.16--1.26$\times$ and 1.39--1.52$\times$. % CLAIMS F4, pinned (paced cells)
The mapped mode costs the host 24.2--26.5\% less than named capture with the copying sender
and 19--27\% less with an object sender, paced, and 35\% less in the fast full-speed regime.
The whole difference is the receiving process, which builds no objects and frees no pages;
outside the process the two are within 0.005 CPU-seconds per gigabyte. Binding the named store
to the mapped-mode area does not close that gap, since the per-page object work remains: bound
capture costs within 2\% of unbound, and the discarding sink is cheaper than all of them.
% CLAIMS F9, pinned
A later series of four commits to capture, in the evaluation kernel only, which releases
backings by kind, allocates extent chunks in blocks, holds pool pages by a reference on the
pool and takes over the fragments of a consumed packet, cuts named capture's full-speed cost by
11\% and bound capture's by 14\% in
the fast regime, with the three controls moving between $-$0.1\% and $+$5.8\%, and narrows the
mapped mode's lead there from 35\% to 23\%; paced, named capture moves within the controls'
spread and the lead is 18.0--27.5\%. % CLAIMS F12, pinned, before the review fixes
% The cheaper-capture series is a follow-up, not in the RFC (docs/design_spec.md §7(11) as it
% stood on 2026-11-07: git show eaec3af:docs/design_spec.md).
Full-speed receive on these machines falls into three regimes per run, two of them separated
by three times the receive interrupts per megabyte, where every receiver's cost outside the
process doubles; only the fast regime is compared. With the driver change of
\S\ref{sec:mapped}, the mode registers on the same interfaces on which the hardware mode is
refused, its streams verified byte for byte with every byte by
reference, and on the kernel with the area-release fix every one of 244 runs completed,
registrations included (\S\ref{sec:mapped}). % CLAIMS F6, F9
\todom{the regimes' cause; the copying NODEV mode; named capture's cost per chunk, which the
receive-side profile attributes to per-fragment accounting and page references and which is
not yet a claim row}

\paragraph{Where hardware zero-copy receive runs.} \todom{on a header-split NIC: the relay at
parity with zero-copy receive plus zero-copy send, and the holding, fan-out and density cases
where the mapped designs copy or starve (W2, Q2)}

\subsection{Zero copy without overhead}
\label{sec:eval:tx}

The send side compares every mechanism on one host sending one stream.
Figure~\ref{fig:tx} plots, and Table~\ref{tab:tx} gives, CPU per gigabyte on the machine under
test for the senders, paced and at full speed, from 1\KiB to 1\MiB: a copying \code{send()}; \sendzc with an unregistered buffer, pinned per send,
in linked chains of eight; \sendzc with a registered buffer in the same chains; the registered
vectored \sendzc, eight buffers in one request, with one request in flight and in chains; the
named mode with one object per send and eight in flight; and the named vectored send, eight
ranges in one request, with one and with eight in flight. Every \sendzc sender bounds its
notification chains (trap 6). The registered \sendzc never rewrites its buffer, so it never
waits for a notification: it is zero copy's best case, and in all but name it is kernel-owned
memory.

% Panel of the final-build TCP comparison. Populate from the qualified campaign.
\begin{minipage}[t]{0.315\textwidth}
\centering
\figplaceholder{2.20cm}{\textbf{(b) Send}\\
Copy, per-send pin, registered pool\\
Named and mapped objects\\[4pt]
\emph{Results pending}}
\end{minipage}


\begin{table*}[t]
\centering
\small
\setlength{\tabcolsep}{4pt}
\begin{tabular}{@{}llrrrrrrrr@{}}
\toprule
rate & size & \code{send()} & \sendzc pinned & \sendzc reg. & reg. vec8 d1 & reg. vec8 d8 & named & named vec8 d1 & named vec8 d8 \\
\midrule
0.5\GBs & 1\KiB & \textbf{0.700} & 1.513 & 1.349 & 0.979 & 0.968 & 1.368 & 1.069 & 1.019 \\
0.5\GBs & 4\KiB & 0.409 & 0.413 & 0.363 & 0.296 & 0.287 & 0.380 & 0.302 & \textbf{0.279} \\
1\GBs & 16\KiB & 0.344 & 0.284 & 0.271 & 0.250 & 0.259 & 0.266 & 0.250 & \textbf{0.237} \\
1\GBs & 64\KiB & 0.316 & 0.268 & 0.260 & 0.248 & 0.238 & 0.232 & 0.237 & \textbf{0.225} \\
full & 1\KiB & 1.300 & -- & -- & 0.367 & \textbf{0.337} & 1.348 & 0.412 & 0.379 \\
full & 4\KiB & 0.293 & 0.360 & 0.190 & 0.106 & \textbf{0.097} & 0.308 & 0.113 & 0.104 \\
full & 16\KiB & 0.170 & 0.130 & 0.104 & 0.083 & 0.090 & 0.102 & 0.074 & \textbf{0.071} \\
full & 64\KiB & 0.158 & 0.101 & 0.083 & 0.073 & 0.067 & 0.076 & 0.067 & \textbf{0.066} \\
full & 256\KiB & 0.145 & 0.084 & 0.072 & \textbf{0.064} & 0.067 & 0.066 & 0.065 & 0.066 \\
full & 1\MiB & 0.148 & 0.087 & 0.065 & 0.063 & 0.063 & 0.063 & 0.066 & \textbf{0.061} \\
\bottomrule
\end{tabular}
\caption{Sender CPU in CPU-seconds per gigabyte delivered, whole host, interrupts pinned, three
runs per cell, every \sendzc sender with bounded notification chains; the cheapest sender per
row in bold. ``d1'' and ``d8'' are one and eight requests in flight; ``vec8'' is eight buffers
or ranges per request. The two non-vectored \sendzc senders at 1\KiB full speed are left out:
the comparison is unresolved there, the pinned sender's nine runs, with bounded chains, a
larger completion queue and unbounded, spreading from 1.27 to 4.61. \msgzc and \sendfile are
compared in the text. \todom{the mapped mode's pages sent through the binding}}
\label{tab:tx}
\end{table*}

\paragraph{The pin.} Against \sendzc that pins per send, like for like, one object per send
against one buffer per send, the named mode is 6--27\% cheaper: at full speed 15\% at 4\KiB,
22\% at 16\KiB, 24\% at 64\KiB, 22\% at 256\KiB and 27\% at 1\MiB; paced, 10\%, 8\%, 6\% and
13\% from 1 to 64\KiB. At 1\KiB full speed the comparison is unresolved: the named mode's runs
lie between 1.22 and 1.43 CPU-seconds per gigabyte, and the pinned sender's nine runs, with
bounded chains, a larger completion queue and unbounded, between 1.27 and 4.61.
% CLAIMS S1, pinned
Pinning itself costs \sendzc 17--90\% at full speed against its own registered buffer from
4\KiB, 90\% at 4\KiB, 25\% at 16\KiB, 22\% at 64\KiB, 17\% at 256\KiB and 33\% at 1\MiB, and
3--14\% paced, or 17\% at 1\KiB and 3\% at 16\KiB without the slow runs (trap 7).
% CLAIMS S2, pinned

\paragraph{The notification.} \sendzc and \msgzc senders must keep a buffer unchanged until its
notification arrives, and an application that reuses buffers waits for it; the named mode's
objects are immutable kernel data and need no reuse tracking. By design, no sender here
measures that wait, since none rewrites its buffer. % CLAIMS S10, argued
\msgzc, the blocking \code{send()} with the zero-copy flag and one buffer per call, pays the
pin, the notification through the error queue and a system call per buffer. Against it the named
mode is 7--48\% cheaper like for like in buffer size: at full speed 24\% at 1\KiB, 31\% at
4\KiB, 48\% at 16\KiB, 46\% at 64\KiB, 33\% at 256\KiB and 12\% at 1\MiB, and 7--14\%
paced. % CLAIMS S11, pinned

\paragraph{Against the best case.} Against the registered, never-rewritten buffer, batched like
for like, eight ranges against eight buffers per request, the named mode ties from 16\KiB and
when paced: at full speed from 16\KiB between 21\% cheaper and 4\% dearer, with the 21\% at
16\KiB and eight in flight, where \sendzc's three runs spread from 0.076 to 0.106; paced,
between 8\% cheaper and 9\% dearer over all runs, and between 11\% cheaper and 4.5\% dearer
without the runs in the sender's slow state (trap 7). At 1--4\KiB saturation it costs 6--12\%
more: 12\% at 1\KiB with one request in flight and with eight, 6\% and 8\% at 4\KiB.
% CLAIMS S3, pinned
Two causes suggested themselves, both named in \S\ref{sec:named:send}: a page reference per fragment that
the registered buffer's managed fragments avoid, and a competitor that sleeps. \sendzc charges a
full page of send buffer for every sub-page fragment, so at 1\KiB only 23\% of its charged bytes
are data, against 93\% for the named mode; its send queue stops near 1\,MB at saturation and it
sleeps about 31\% of the time, which keeps it off the socket lock. % CLAIMS S9, measured
Doing the named mode's per-send bookkeeping with the socket owned narrowed the gap within one
campaign from 28--29\% to 10--18\% at 1\KiB and from 20--30\% to 12.5--20\% at 4\KiB; the lock
holders are TCP's own deferred transmits, which can only start while the socket is unowned.
% CLAIMS S7 (pinned), S8 (in-placement probes)
Changing TCP's autocork policy for fragment-limited streams, in two experiment kernels, moved
neither sender beyond its run-to-run spread and did not narrow the gap; with \sendzc's chains
bounded, the gap is the 6--12\% above. % CLAIMS M4, pinned

\paragraph{The special case.} With the packet tap of \S\ref{sec:eval:method} in the transmit
path, \sendzc's pages are copied and the named mode's are cloned: the tap costs \sendzc
15\%, 32\% and 41\% paced and 60--149\% at full speed, and the named mode 1--5\%
(\S\ref{sec:eval:method}, trap 2). % CLAIMS F7, pinned
A \sendzc request can opt out of that copy on a development branch of \iouring that manages
its fragments; a tap then copies none of its packets and costs it 2.5\% at 1\KiB, 2\% at 4\KiB
and 3\% at 64\KiB paced and between 3\% less and 9\% more at full speed, and with the tap out
of the path it costs within 4\% of plain \sendzc. Managed fragments for
the named mode's sends, one folio reference per packet instead of a page reference per
fragment, save 0--5\%, and taps still do not copy them; what is left of the 1--4\KiB gap to
the registered buffer is the sending process's own CPU, not the fragments.
% CLAIMS S13, S14, pinned
On contiguous ranges the named mode's fragments coalesce and the registered buffer's
managed fragments never do: eight consecutive 1\KiB ranges per request cost the named mode 60\%
less than \sendzc at 0.5\GBs and 11\% less at full speed; at 4\KiB it costs 5\% less paced and
the registered buffer is 11\% cheaper at full speed. % CLAIMS S4, pinned

\paragraph{What remains.} Profiles of the sending thread put the rest of the small-send gap in
the named mode's request path, not in its model. Fixing the managed fill cut the fill's own
work but moved no end-to-end number, and prepaying folio references per send moves the
per-packet stall instead of removing it, so prepaying is dropped.
% CLAIMS S15, pinned; prepay removed from the development branch (ledger 42)
Fixing the request's issue path instead, admission in one pass, a checked copy-in loop and one
message-header zeroing per issue, cut the request from 54 to 37 CPU-milliseconds per gigabyte
at 1\KiB full speed and from 63 to 43 paced, to within 5 CPU-milliseconds per gigabyte of
\sendzc's request at full speed and within 1 at 4\KiB; end to end that is about 2\% of the
host's CPU at 1\KiB paced, below what total minus idle resolves at four runs per kernel.
% CLAIMS S17, pinned (the commits are on the development branch, without prepay)
A further set of commits on the send's request state changes no paced cell beyond the
run-to-run spread. % CLAIMS S18, pinned (development branch, without prepay)
What is left of the single-object operation, 132--145\,ns per send, is mostly decisions of the
design: the request count's two locked instructions, the object table's mutex, the unpin and
the ordering come to 62--70\,ns, against 30--33\,ns for \sendzc's operation and 86--98\,ns for
its notification and completion work. The unpin is an atomic only because these senders send
the same objects again; on a relay's last send of an object, which detaches it, there is no
atomic to remove. % CLAIMS S19, pinned (the unpin note: ledger 43)
Removing the request count with its limit and making the table's lock a spinlock since took
14--17\,ns per send off the operation at 1\KiB, of the 20--25 estimated; part of it reappears
in \iouring's core and in TCP's transmit path, and end to end the named mode moved between
$-$3.4\% and 0 in the paced cells while the controls moved between $-$4.9\% and $+$2.7\%.
% CLAIMS S20, pinned
All of these changes are in the evaluation kernel only; the series of \S\ref{sec:impl} has no
request count either, but keeps the table's mutex. % design_spec.md §2.2, §2.7, §7(8), §7(11) as of 2026-11-07 (git show eaec3af:docs/design_spec.md)

\paragraph{Vectoring and the break-even.} Eight ranges in one request make the named mode
3.6$\times$ cheaper at 1\KiB full speed and 3.0$\times$ at 4\KiB than one object per send, and
1.3--1.4$\times$ paced; the copying \code{send()} gains the same way from one request per eight
buffers. % CLAIMS S5, pinned
Against copying, the named mode is 40--57\% cheaper from 16\KiB at full speed and its vectored
send 62--71\% cheaper at 1--4\KiB. But at 0.5\GBs with 1\KiB messages, copying is the cheapest
sender of all, 0.700 against 0.968--1.513 for every zero-copy sender, and one object per send
costs about what \code{send()} does at 1 and 4\KiB full speed, 4\% and 5\% more, within the
named mode's spread: every zero-copy path pays a fixed cost per operation, and below the
break-even the copy of a kilobyte is cheaper than the cheapest of them. % CLAIMS S6, pinned

\paragraph{Against \sendfile.} The page cache with \sendfile is the kernel's own way to send held
data by reference, and the baseline that a cache built on files would use. Sending one object
per request, the named mode costs 28--39\% less than \sendfile from a page cache of 4\KiB
folios at full speed from 64\KiB to 1\MiB, and 10\% less at 1\GBs. % CLAIMS S12, pinned
An earlier version of the splice series of \S\ref{sec:impl} took a 64\KiB packet from 16
fragments to 2, for \sendfile from large folios and for object sends alike. It made object
sends 24\% cheaper at 64\KiB full speed, 14\% at 256\KiB and 5\% at 1\GBs, and \sendfile
from 64\KiB folios 11\% cheaper at full speed and 5\% at 1\GBs, while \sendfile from
4\KiB folios did not change; the series as prepared has not been measured.
% CLAIMS N1, N2, N3 (an earlier implementation; design_spec.md §4.3)

\subsection{Forwarding}
\label{sec:eval:fwd}

The relay of \S\ref{sec:apps} forwards one stream through the machine under test, with plain
endpoints on the other two machines, at 1\GBs and at full speed.

% Corrected 2026-11-11: this form said that the named relay ran with inline compaction and copied
% every captured byte once, that the paced cost was that copy, and that binding removed it. The
% relay sends no COMPACT request and copies nothing (CLAIMS F1).
\paragraph{Parity with in-kernel forwarding.} The named relay forwards by reference and copies
nothing. It costs
11\% more CPU per gigabyte than \splice at 0.5 and 1\GBs, and 4\% more at full speed, where
it delivers 21\% more, 2.38 against 1.96\GBs; against the copying relay it costs 2--5\% more
paced and 23\% less at full speed. % CLAIMS F1, pinned; the 21% is not attributed to the relay

\paragraph{Forwarding from lent memory.} Binding the relay's store to a mapped-mode area keeps
captured bytes in the area, and the relay then costs
14\% less CPU per gigabyte than \splice at 0.5\GBs, 15\% less at 1\GBs and 16\% less at
full speed at equal throughput, which is 22--24\% less than the unbound named relay. A plain
mapped-mode relay, zero-copy receive into \sendzc with a pin per send and each buffer returned
after its notification, costs 15\% and 17\% less than
\splice paced and 30\% less at full speed while delivering 11\% less. With the sink on the
mapped mode as well, the relay costs 19\% less at 1\GBs and 22\% less at full speed than
with both ends named. Small transactions do not gain: bound, the relay's median is
150\,$\mu$s against 142 for \splice, and it spends 19\% more CPU per transaction, as the
unbound relay does. % CLAIMS F10, pinned

\paragraph{Request/response.} Sixty-four-byte transactions see a median of 59\,$\mu$s
direct and 146--159\,$\mu$s through any of the forwarders: the network hop sets the latency,
not the forwarder. Per million transactions the machine under test spends 30.8 CPU-seconds
with the copying relay, 32.9 with \splice, 28.9 with sockmap and 37.0 with the named relay,
which is 12--20\% more than \splice or the copying relay: the price of a name, paid per small
transaction (\S\ref{sec:lending:cost}). % CLAIMS F8, pinned

\paragraph{What the in-kernel alternatives give up.} Sockmap's redirect has no backpressure: a
verdict program appends each packet to the egress socket's queue with no length or memory
limit, so the ingress connection is acknowledged as fast as the verdict runs and the queue
grows in kernel memory, invisible from userspace and uncharged to TCP's memory accounting.
At full speed it delivered 0.64--1.08\GBs with the queue growing in every run, and paced it was
the dearest zero-copy forwarder, 11--50\% over the named relay; with the interrupts pinned it
is the dearest forwarder at every rate and delivers 0.81\GBs at full speed with an erratic
cost. % CLAIMS F2, measured; F1, pinned
A correct teardown needs a workaround: the only way to know the queue is drained is to wait for
the egress socket's acknowledged byte count to cover every redirected byte.

\paragraph{Compaction.} Compaction with a cold destination, cut at every fragment and run from an
unbound worker, cost 2.1$\times$ a relay that copies every byte twice; the redesign of
\S\ref{sec:named:memory}, one copy per contiguous run into recycled store-owned units on the
forwarder's own CPU, cut it by 26--33\%, and its round-trip latency fell to that of the named
relay without compaction. % CLAIMS F3, measured
\todom{compaction's remaining cost over the plain named relay on the final kernel, and its
effect on data-path latency while it runs (H2)}

\paragraph{Keep, mirror, replay.} \todom{the capability demonstration: a relay that keeps the
body, mirrors it to a second backend and replays it after a backend failure, against a copying
proxy that buffers for retry; \splice cannot, \code{tee} mirrors within a pipe (X4)}

\subsection{Retention}
\label{sec:eval:retention}

\todom{holding at equal memory: bytes pinned per payload byte by driver and value size, with
compaction off and on, and capacity at equal physical memory against the copying store (H1);
fan-out: memory and CPU against $N$ receivers for one object, against a copy per receiver, a
RAM file with \sendfile and \splice with \code{tee} (X3, H4); density: applications per NIC,
unlimited in the named mode, one per device in the mapped mode and a queue each for zero-copy
receive (W3); behaviour at the budget: slow and hostile readers, holding past the budget,
teardown (H3)}

\subsection{The store}
\label{sec:eval:store}

\todom{the store with its data path switched between copy, zero-copy send, a RAM file with
\sendfile, \splice with \code{tee}, zero-copy receive and the named mode, under equal memory
budgets; then against Redis, Valkey and memcached from tens of bytes to megabytes: CPU per
request and per gigabyte against value size in one figure, throughput, tail latency, capacity
at equal memory (X2, X5)}
